% =====================================================================
%  Figura: pipeline di trasmissione della boa (CSMA/CA)
%
%  Come usarla in un capitolo:
%     % =====================================================================
%  Figura: pipeline di trasmissione della boa (CSMA/CA)
%
%  Come usarla in un capitolo:
%     % =====================================================================
%  Figura: pipeline di trasmissione della boa (CSMA/CA)
%
%  Come usarla in un capitolo:
%     % =====================================================================
%  Figura: pipeline di trasmissione della boa (CSMA/CA)
%
%  Come usarla in un capitolo:
%     \input{fig_tx_pipeline_it}
%  dove vuoi che compaia la figura (ad es. nel Cap. 3, vicino a
%  \label{sec:fsm}). Richiede solo \usepackage{graphicx}, gia' presente
%  nel preambolo.
%
%  L'immagine e' generata da thesis/make_pipeline_figure.py e salvata in
%  thesis/figures/tx_pipeline.pdf (vettoriale) + .png. Per rigenerarla:
%     .venv/Scripts/python.exe thesis/make_pipeline_figure.py
% =====================================================================
\begin{figure}[p]
  \centering
  \includegraphics[height=0.86\textheight,keepaspectratio]{figures/tx_pipeline}
  \caption[Pipeline di trasmissione della boa]{%
    Pipeline di trasmissione di una boa (CSMA/CA), come implementata in
    \texttt{buoy.py}. Lo scheduler decide periodicamente se trasmettere; in
    caso affermativo il nodo avvia un'unica pipeline protetta dal lock di
    \texttt{processing}, che rileva il canale prima e dopo il DIFS e al
    termine del backoff (Eq.~\ref{eq:backoff}): se il canale risulta occupato
    in una qualunque delle tre verifiche, si rientra dal \textsc{channel\_sense}.
    Vinta la contesa, la boa invia un solo frame -- con precedenza al beacon
    proprio, altrimenti un inoltro dalla coda con il TTL decrementato
    (Eq.~\ref{eq:ttl}) -- e, se restano inoltri in sospeso, rientra subito in
    contesa. Per leggibilit\`a non \`e mostrato il controllo, ripetuto a ogni
    stadio, che rilascia il lock se non resta nulla da inviare.}
  \label{fig:tx-pipeline}
\end{figure}

%  dove vuoi che compaia la figura (ad es. nel Cap. 3, vicino a
%  \label{sec:fsm}). Richiede solo \usepackage{graphicx}, gia' presente
%  nel preambolo.
%
%  L'immagine e' generata da thesis/make_pipeline_figure.py e salvata in
%  thesis/figures/tx_pipeline.pdf (vettoriale) + .png. Per rigenerarla:
%     .venv/Scripts/python.exe thesis/make_pipeline_figure.py
% =====================================================================
\begin{figure}[p]
  \centering
  \includegraphics[height=0.86\textheight,keepaspectratio]{figures/tx_pipeline}
  \caption[Pipeline di trasmissione della boa]{%
    Pipeline di trasmissione di una boa (CSMA/CA), come implementata in
    \texttt{buoy.py}. Lo scheduler decide periodicamente se trasmettere; in
    caso affermativo il nodo avvia un'unica pipeline protetta dal lock di
    \texttt{processing}, che rileva il canale prima e dopo il DIFS e al
    termine del backoff (Eq.~\ref{eq:backoff}): se il canale risulta occupato
    in una qualunque delle tre verifiche, si rientra dal \textsc{channel\_sense}.
    Vinta la contesa, la boa invia un solo frame -- con precedenza al beacon
    proprio, altrimenti un inoltro dalla coda con il TTL decrementato
    (Eq.~\ref{eq:ttl}) -- e, se restano inoltri in sospeso, rientra subito in
    contesa. Per leggibilit\`a non \`e mostrato il controllo, ripetuto a ogni
    stadio, che rilascia il lock se non resta nulla da inviare.}
  \label{fig:tx-pipeline}
\end{figure}

%  dove vuoi che compaia la figura (ad es. nel Cap. 3, vicino a
%  \label{sec:fsm}). Richiede solo \usepackage{graphicx}, gia' presente
%  nel preambolo.
%
%  L'immagine e' generata da thesis/make_pipeline_figure.py e salvata in
%  thesis/figures/tx_pipeline.pdf (vettoriale) + .png. Per rigenerarla:
%     .venv/Scripts/python.exe thesis/make_pipeline_figure.py
% =====================================================================
\begin{figure}[p]
  \centering
  \includegraphics[height=0.86\textheight,keepaspectratio]{figures/tx_pipeline}
  \caption[Pipeline di trasmissione della boa]{%
    Pipeline di trasmissione di una boa (CSMA/CA), come implementata in
    \texttt{buoy.py}. Lo scheduler decide periodicamente se trasmettere; in
    caso affermativo il nodo avvia un'unica pipeline protetta dal lock di
    \texttt{processing}, che rileva il canale prima e dopo il DIFS e al
    termine del backoff (Eq.~\ref{eq:backoff}): se il canale risulta occupato
    in una qualunque delle tre verifiche, si rientra dal \textsc{channel\_sense}.
    Vinta la contesa, la boa invia un solo frame -- con precedenza al beacon
    proprio, altrimenti un inoltro dalla coda con il TTL decrementato
    (Eq.~\ref{eq:ttl}) -- e, se restano inoltri in sospeso, rientra subito in
    contesa. Per leggibilit\`a non \`e mostrato il controllo, ripetuto a ogni
    stadio, che rilascia il lock se non resta nulla da inviare.}
  \label{fig:tx-pipeline}
\end{figure}

%  dove vuoi che compaia la figura (ad es. nel Cap. 3, vicino a
%  \label{sec:fsm}). Richiede solo \usepackage{graphicx}, gia' presente
%  nel preambolo.
%
%  L'immagine e' generata da thesis/make_pipeline_figure.py e salvata in
%  thesis/figures/tx_pipeline.pdf (vettoriale) + .png. Per rigenerarla:
%     .venv/Scripts/python.exe thesis/make_pipeline_figure.py
% =====================================================================
\begin{figure}[p]
  \centering
  \includegraphics[height=0.86\textheight,keepaspectratio]{figures/tx_pipeline}
  \caption[Pipeline di trasmissione della boa]{%
    Pipeline di trasmissione di una boa (CSMA/CA), come implementata in
    \texttt{buoy.py}. Lo scheduler decide periodicamente se trasmettere; in
    caso affermativo il nodo avvia un'unica pipeline protetta dal lock di
    \texttt{processing}, che rileva il canale prima e dopo il DIFS e al
    termine del backoff (Eq.~\ref{eq:backoff}): se il canale risulta occupato
    in una qualunque delle tre verifiche, si rientra dal \textsc{channel\_sense}.
    Vinta la contesa, la boa invia un solo frame -- con precedenza al beacon
    proprio, altrimenti un inoltro dalla coda con il TTL decrementato
    (Eq.~\ref{eq:ttl}) -- e, se restano inoltri in sospeso, rientra subito in
    contesa. Per leggibilit\`a non \`e mostrato il controllo, ripetuto a ogni
    stadio, che rilascia il lock se non resta nulla da inviare.}
  \label{fig:tx-pipeline}
\end{figure}
